\environment env_ConfigGlobal
\environment env_Config

\startcomponent Exercises1
\startsection[title={Exercises for Lesson 1}, reference={sec:exercises1}]

Let's try making sentences using what we have learned so far. You can use romaji,
hiragana, or full Japanese (hiragana, katakana and any kanji you know).
Vocabulary is listed at the end.

\startsubsubject[title={Sentences}]

\startitemize[n]
\item I am a cat
\item I dance
\item The cat walks
\item The bird sings
\item The cat is white
\item The bird is small
\stopitemize

\stopsubsubject

\startsubsubject[title={Vocabulary}]

\startitemize[packed]
\item I: \JPN{わたし} watashi
\item Cat: \JPN{ねこ} neko
\item Dance: \JPN{おどる} odoru
\item Walk: \JPN{あるく} aruku
\item Bird: \JPN{とり} tori
\item Sing: \JPN{うたう} utau
\item Is-white: \JPN{\ruby{白}{しろ}い} shiroi
\item Is-small: \JPN{\ruby{小}{ちい}さい} chiisai
\stopitemize

Here I am using only fairly basic kanji---broadly those covered in {\em Alice in
Kanji Land}. Feel free to use more kanji if you know them. I am giving animals,
birds and plants in hiragana, but katakana is often used for them and you may do
so if you prefer (not {\em tori\/} though---the actual names of particular animals,
birds, and plants such as \JPN{ウサギ} or \JPN{ネコ}). Using katakana is probably
slightly more \quotation{correct}-seeming but hiragana is considered softer and
therefore cuter. I am also doing it to make it easier for those who only know
hiragana. We should work in (at least) hiragana from an early stage for the sake
of structural integrity. Taking longer with katakana isn't a crucial issue.

\stopsubsubject

\stopsection
\stopcomponent
